% =====================================================================
%  Matrix Methods for Optimizing Industrial Piping Networks
%  Adjacency Matrices, Graph Laplacians, and Cost Flow Analysis
%  Applied to Oil Refinery Design
%
%  Course:   Matrices and Linear Transformations (BCSAI, 2026)
%  Status:   FULL DRAFT (all 8 sections present).
%
%  This file is standalone so it compiles on its own (Overleaf or
%  local TeX Live).  It is shared through the group git repo because
%  Overleaf free only allows a single editor at a time.
%
%  Sections in this file (in document order):
%    - Title
%    - Abstract
%    - 1. Introduction
%    - 2. Theory and Mathematical Formulas
%    - 3. Research and Application (6-node refinery case study)
%    - 4. Conclusion
%    - Bibliography (textbook seeds; extend with engineering sources)
%    - Annex A (Python code for Layouts A and B)
%
%  Left for teammates:
%    - Review and polish of the draft text; adding figures where
%      "% FIGURE:" placeholders appear.
%    - Expanding the Bibliography beyond the six textbook seeds.
%    - Building the presentation slides (a separate deliverable).
% =====================================================================

\documentclass[11pt,a4paper]{article}

% ---- Packages --------------------------------------------------------
\usepackage[a4paper,margin=1in]{geometry}
\usepackage{amsmath,amssymb,amsthm}
\usepackage{graphicx}
\usepackage{booktabs}
\usepackage{xcolor}
\usepackage{cite}
\usepackage[hidelinks]{hyperref}

% ---- Shortcuts for consistent notation ------------------------------
\newcommand{\RR}{\mathbb{R}}
\newcommand{\mat}[1]{\mathbf{#1}}
\newcommand{\vect}[1]{\mathbf{#1}}
\newcommand{\transpose}{^{\mathsf{T}}}

% ---- Title -----------------------------------------------------------
\title{Matrix Methods for Optimizing Industrial Piping Networks:\\
Adjacency Matrices, Graph Laplacians, and Cost Flow Analysis\\
Applied to Oil Refinery Design}
\author{Group Project \\ Matrices and Linear Transformations (BCSAI, 2026)}
\date{}

\begin{document}
\maketitle

% =====================================================================
%  ABSTRACT
% =====================================================================
\begin{abstract}
Oil refineries are, at heart, very large piping networks, and the
decisions that shape them, which units to connect, which pipes to
duplicate, how much to spend on each link, are naturally linear
algebra decisions. This report models a simplified refinery as a
directed graph with six nodes (a crude intake, two distillation
towers, an internal junction, and two product outlets) and seven
pipes, and uses only the three topics assigned to our group,
\emph{Matrices}, \emph{Gauss Elimination}, and \emph{Eigenvalues}, to
answer three questions about it. First, we write mass conservation
as the linear system $\mat{B}\vect{f}=\vect{s}$, reduce it by
Gaussian elimination, and describe the full two-parameter family of
feasible flows. Second, we compute the eigenvalues of the graph
Laplacian analytically and use the Fiedler value $\lambda_2 =
3-\sqrt{5} \approx 0.764$ to quantify how robust the layout is to
the loss of a pipe. Third, we use a linear cost model to compare two
candidate layouts and show that the reference layout is cheaper (cost
$400$ vs.\ $460$) \emph{and} structurally more robust than a layout
missing one pipe. The same matrix--eigenvalue--cost workflow scales
directly to realistic refinery design.
\end{abstract}

% =====================================================================
%  1. INTRODUCTION  (target ~1.5 pages)
% =====================================================================
\section{Introduction}
\label{sec:introduction}

An oil refinery is, at its core, a very large and very expensive network
of pipes. Crude oil enters through a handful of intake lines, is split
between distillation columns, reactors, heat exchangers, and storage
tanks, and eventually leaves as finished products such as gasoline,
diesel, jet fuel, or petrochemical feedstocks. Between any two of these
units sits a pipe, and between any two pipes sits a junction, a pump, or
a valve. \emph{Optimizing} such a network means choosing a layout that
respects three kinds of constraints simultaneously: the physical
conservation of mass at every junction, the capacity limits of each
pipe, and the cost of actually building and operating the network.
Because industrial piping decisions are made once and live for decades,
getting this right matters: a bad layout either bottlenecks throughput,
raises pumping costs, or worse, creates a single point of failure that
can shut the whole refinery down.

The key observation of this report is that all three of the constraints
above are \emph{linear}. Flow balance at a junction is a linear
equation. Pipe capacity is a linear inequality. Total operating cost is
a linear function of the flows carried by each pipe. The network itself
is a graph, and every graph can be represented by matrices: the
adjacency matrix records which nodes are connected, the incidence matrix
translates pipe flows into node-level balances, and the graph Laplacian
encodes how well the network holds together as a whole. Solving the
resulting system of equations is exactly what Gaussian elimination was
designed for, and the eigenvalues of the Laplacian are exactly what tell
us whether the network is well-connected or fragile. In other words,
the three topics assigned to our group, \emph{matrices}, \emph{Gauss
elimination}, and \emph{eigenvalues}, are not only relevant to refinery
design, they are the natural language for it.

This project applies that language to a concrete design question. We
model a simplified oil refinery piping network as a directed graph,
build its adjacency, degree, incidence, and Laplacian matrices, and use
Gaussian elimination to solve the resulting flow equations. We then
examine the eigenvalues of the Laplacian to identify bottlenecks and
measure the network's robustness to the loss of a pipe, and we use a
linear cost model to compare alternative layouts. The aim is not to
design a working refinery in fifteen pages but to demonstrate, with
small but honest examples, that the same linear algebra we learn in
class is the backbone of real industrial optimization.

The rest of the document is organized as follows. Section
\ref{sec:theory} develops the theory and fixes a notation that the
later sections will reuse without change: graphs and their
representation matrices, the Graph Laplacian and its spectrum, linear
systems for conservation laws, Gaussian elimination, eigenvalues, and
a linear cost--flow model. The Research and Application section that
follows applies this framework to a specific refinery example. The
Conclusion summarizes the findings and points to avenues for further
work, and the Bibliography lists the textbooks and references that
support the arguments throughout.

% =====================================================================
%  2. THEORY AND MATHEMATICAL FORMULAS  (target ~3 pages)
% =====================================================================
\section{Theory and Mathematical Formulas}
\label{sec:theory}

Before presenting the mathematics, we fix a notation that will be used
throughout the report. Every symbol introduced here keeps the same
meaning in later sections; the Research and Application section can
plug in refinery-specific values without redefining anything.

\paragraph{Notation.}
\begin{itemize}
  \item $G = (V, E)$ denotes a graph with $n = |V|$ nodes and
        $m = |E|$ edges. In the refinery context, nodes are
        junctions, pumps, tanks, and processing units, and edges are
        pipes.
  \item $\mat{A} \in \RR^{n \times n}$ is the \emph{adjacency matrix}.
  \item $\mat{D} \in \RR^{n \times n}$ is the \emph{degree matrix}.
  \item $\mat{B} \in \RR^{n \times m}$ is the \emph{incidence matrix}.
  \item $\mat{L} = \mat{D} - \mat{A} \in \RR^{n \times n}$ is the
        \emph{Graph Laplacian}.
  \item $\vect{f} \in \RR^{m}$ is a \emph{flow vector}: one entry per
        pipe, giving the amount of material flowing along that pipe in
        its chosen direction.
  \item $\vect{s} \in \RR^{n}$ is an \emph{external supply vector}:
        $s_i > 0$ when node $i$ injects material into the network (a
        source such as a crude intake), $s_i < 0$ when it removes
        material (a sink such as a product outlet), and $s_i = 0$ at
        interior junctions.
  \item $\vect{c} \in \RR^{m}$ is a \emph{cost vector}, with $c_e$ the
        unit cost of pushing flow through pipe $e$.
\end{itemize}

% --------------------------------------------------------------------
\subsection{Graphs and the matrices that represent them}
\label{sec:graphs}

A piping network is naturally a graph. A \emph{graph} $G = (V, E)$ is a
pair of sets: $V$ is a finite set of \emph{nodes}, and $E$ is a set of
\emph{edges} connecting pairs of nodes. An \emph{undirected} edge is an
unordered pair $\{u, v\}$; a \emph{directed} edge is an ordered pair
$(u, v)$. For flow problems we use the directed version, since a pipe
carries material in a chosen direction.

\paragraph{Adjacency matrix.}
The adjacency matrix $\mat{A} \in \RR^{n \times n}$ records which nodes
share an edge:
\[
A_{ij} =
\begin{cases}
  1 & \text{if there is an edge between node $i$ and node $j$,} \\
  0 & \text{otherwise.}
\end{cases}
\]
For an undirected graph, $\mat{A}$ is symmetric.

\paragraph{Degree matrix.}
The \emph{degree} $d_i$ of node $i$ is the number of edges incident to
$i$. The degree matrix is the diagonal matrix
$\mat{D} = \operatorname{diag}(d_1, d_2, \ldots, d_n)$.

\paragraph{Incidence matrix.}
For a directed graph with $m$ edges, the incidence matrix
$\mat{B} \in \RR^{n \times m}$ has one row per node and one column per
edge. For each directed edge $e$ going from node $u$ to node $v$, set
\[
B_{u,e} = +1, \qquad B_{v,e} = -1,
\]
and all other entries in column $e$ equal to zero. With this sign
convention, the product $\mat{B}\vect{f}$ at row $i$ equals the
\emph{net outflow} at node $i$: the sum of flows on pipes leaving node
$i$ minus the sum of flows on pipes entering it. This is the
convention we will use everywhere in this document.

\paragraph{A small worked example.}
Consider the four-node graph with edges
$e_1 = (1 \to 2)$, $e_2 = (2 \to 3)$, $e_3 = (3 \to 4)$, and
$e_4 = (1 \to 3)$. Its adjacency, degree, and incidence matrices are
\[
\mat{A} =
\begin{bmatrix}
  0 & 1 & 1 & 0 \\
  1 & 0 & 1 & 0 \\
  1 & 1 & 0 & 1 \\
  0 & 0 & 1 & 0
\end{bmatrix},
\quad
\mat{D} =
\begin{bmatrix}
  2 & 0 & 0 & 0 \\
  0 & 2 & 0 & 0 \\
  0 & 0 & 3 & 0 \\
  0 & 0 & 0 & 1
\end{bmatrix},
\quad
\mat{B} =
\begin{bmatrix}
  \phantom{-}1 & \phantom{-}0 & \phantom{-}0 & \phantom{-}1 \\
  -1 & \phantom{-}1 & \phantom{-}0 & \phantom{-}0 \\
  \phantom{-}0 & -1 & \phantom{-}1 & -1 \\
  \phantom{-}0 & \phantom{-}0 & -1 & \phantom{-}0
\end{bmatrix}.
\]
% FIGURE: teammate to insert a diagram of this 4-node graph (TikZ or
% PNG) here; keep the same labelling (nodes 1..4, edges e1..e4).
The three matrices encode the same network from three complementary
angles: $\mat{A}$ asks ``who is connected to whom'', $\mat{D}$ asks
``how busy is each node'', and $\mat{B}$ asks ``what does a given flow
pattern do at every node''.

% --------------------------------------------------------------------
\subsection{The Graph Laplacian}
\label{sec:laplacian}

Combining the degree and adjacency matrices produces the single most
useful object in network analysis, the \emph{graph Laplacian}:
\begin{equation}
  \mat{L} \;=\; \mat{D} - \mat{A}.
  \label{eq:laplacian}
\end{equation}
For the four-node example of Section~\ref{sec:graphs},
\[
\mat{L} =
\begin{bmatrix}
  \phantom{-}2 & -1 & -1 & \phantom{-}0 \\
  -1 & \phantom{-}2 & -1 & \phantom{-}0 \\
  -1 & -1 & \phantom{-}3 & -1 \\
  \phantom{-}0 & \phantom{-}0 & -1 & \phantom{-}1
\end{bmatrix}.
\]

The Laplacian has four properties that we will use repeatedly.
First, it is \emph{symmetric} whenever the underlying graph is
undirected (or when we symmetrize a directed graph by forgetting
orientation). Second, every row of $\mat{L}$ sums to zero, so the
all-ones vector $\vect{1} = (1, 1, \ldots, 1)\transpose$ is always an
eigenvector with eigenvalue $0$. Third, $\mat{L}$ is \emph{positive
semidefinite}: $\vect{x}\transpose \mat{L} \vect{x} \ge 0$ for every
$\vect{x} \in \RR^{n}$, which means all eigenvalues are real and
non-negative. Fourth, $\mat{L}$ can be written as
$\mat{L} = \mat{B}\mat{B}\transpose$ (for the undirected orientation
convention), which connects the node-level and edge-level views of the
network.

Order the eigenvalues of the Laplacian as
\[
0 = \lambda_1 \le \lambda_2 \le \cdots \le \lambda_n.
\]
The first eigenvalue is always zero, as already noted. The
\emph{second} eigenvalue $\lambda_2$ is called the \emph{Fiedler value}
or \emph{algebraic connectivity}, and it carries a very concrete
meaning: $\lambda_2 > 0$ if and only if the graph is connected, and
the larger $\lambda_2$ is, the harder it is to disconnect the graph
by cutting edges. For a piping network this translates directly:
$\lambda_2$ measures how robust the network is to the loss of a pipe.
A refinery layout with a small $\lambda_2$ has a bottleneck somewhere,
and losing the wrong pipe can break the network into pieces that
cannot exchange material. A layout with a larger $\lambda_2$ has more
redundancy and degrades more gracefully.

% --------------------------------------------------------------------
\subsection{Linear systems and conservation laws}
\label{sec:conservation}

The most important physical law in a piping network is conservation
of mass: whatever flows into a junction must also flow out (except at
sources and sinks, which inject or extract material from the outside).
Using the notation of Section~\ref{sec:graphs}, conservation at every
node can be written as a single linear system:
\begin{equation}
  \mat{B}\vect{f} \;=\; \vect{s}.
  \label{eq:balance}
\end{equation}

Each row of \eqref{eq:balance} is the balance equation for one node.
Reading it from left to right, the $i$-th row says: \emph{net outflow
at node $i$ equals the external supply at node $i$}. At an interior
junction, $s_i = 0$, and \eqref{eq:balance} reduces to ``flow in =
flow out''. At a source node such as a crude oil intake, $s_i > 0$
encodes the rate at which material is injected. At a sink node such
as a product outlet, $s_i < 0$ encodes the extraction rate.

The refinery vocabulary for these objects is worth stating explicitly,
since the Application section will use the same symbols:
\begin{center}
\begin{tabular}{@{}ll@{}}
\toprule
Symbol & Refinery meaning \\
\midrule
$f_e$        & volumetric flow carried by pipe $e$ \\
$s_i > 0$    & net production at node $i$ (e.g. crude intake) \\
$s_i < 0$    & net consumption at node $i$ (e.g. product outlet) \\
$s_i = 0$    & interior junction or passive element \\
$(\mat{B}\vect{f})_i$ & net outflow at node $i$ implied by $\vect{f}$ \\
\bottomrule
\end{tabular}
\end{center}

System \eqref{eq:balance} is typically \emph{underdetermined} for a
realistic refinery: there are more pipes than junctions, so there are
more unknowns than equations, and many flow patterns satisfy mass
balance. Choosing among them is exactly the optimization problem that
the cost model in Section~\ref{sec:cost} will formalize. The method
we use to solve or reduce such systems is the subject of the next
subsection.

% --------------------------------------------------------------------
\subsection{Gaussian elimination}
\label{sec:gauss}

Gaussian elimination is the workhorse method for solving a linear
system $\mat{M}\vect{x} = \vect{b}$. It proceeds by applying elementary
row operations, (i) swapping two rows, (ii) multiplying a row by a
non-zero scalar, and (iii) adding a multiple of one row to another, to
transform the augmented matrix $[\mat{M} \mid \vect{b}]$ into an upper
triangular form. From the triangular system, the unknowns are then
recovered by \emph{back substitution}, starting from the last equation
and working upward. When a pivot (the leading coefficient of the row
currently being used) is zero or very small, we swap in a row with a
larger pivot; this is called \emph{partial pivoting} and is what makes
the method numerically stable in practice.

\paragraph{Worked example.}
Consider the $3 \times 3$ system
\begin{equation}
  \begin{cases}
    \phantom{-}x + \phantom{-}y + \phantom{-}z = \phantom{-}6, \\
    2x + 3y + \phantom{-}z = 11, \\
    \phantom{-}x - \phantom{-}y + 2z = \phantom{-}5.
  \end{cases}
  \label{eq:gauss-sys}
\end{equation}
Write \eqref{eq:gauss-sys} as the augmented matrix
\[
\left[\begin{array}{ccc|c}
1 & \phantom{-}1 & 1 & \phantom{-}6 \\
2 & \phantom{-}3 & 1 & 11 \\
1 & -1 & 2 & \phantom{-}5
\end{array}\right].
\]
Eliminate the first column below the pivot by
$R_2 \gets R_2 - 2R_1$ and $R_3 \gets R_3 - R_1$:
\[
\left[\begin{array}{ccc|c}
1 & \phantom{-}1 & \phantom{-}1 & \phantom{-}6 \\
0 & \phantom{-}1 & -1 & -1 \\
0 & -2 & \phantom{-}1 & -1
\end{array}\right].
\]
Eliminate the second column below the new pivot by
$R_3 \gets R_3 + 2R_2$:
\[
\left[\begin{array}{ccc|c}
1 & 1 & \phantom{-}1 & \phantom{-}6 \\
0 & 1 & -1 & -1 \\
0 & 0 & -1 & -3
\end{array}\right].
\]
The matrix is now upper triangular. Back substitution gives
$z = 3$, then $y - z = -1 \Rightarrow y = 2$, and finally
$x + y + z = 6 \Rightarrow x = 1$, so
$(x, y, z) = (1, 2, 3)$. This is exactly the routine we will apply to
the refinery balance equations in the Application section, except that
the matrix there will come from \eqref{eq:balance}.

Gaussian elimination on an $n \times n$ system costs $\mathcal{O}(n^3)$
arithmetic operations, which is why large-scale networks in industry
rely on sparse variants, iterative solvers, or specialized factorizations
(LU, QR, Cholesky). The textbook references at the end of this document
develop these extensions in detail.

% --------------------------------------------------------------------
\subsection{Eigenvalues and eigenvectors}
\label{sec:eigen}

Given a square matrix $\mat{M} \in \RR^{n \times n}$, a non-zero vector
$\vect{v} \in \RR^{n}$ is an \emph{eigenvector} of $\mat{M}$ with
\emph{eigenvalue} $\lambda \in \RR$ if
\begin{equation}
  \mat{M}\vect{v} = \lambda \vect{v}.
  \label{eq:eigendef}
\end{equation}
Eigenvalues are the roots of the \emph{characteristic polynomial}
$p(\lambda) = \det(\mat{M} - \lambda \mat{I})$, where $\mat{I}$ is the
$n \times n$ identity matrix. They tell us in which directions the
linear map $\vect{x} \mapsto \mat{M}\vect{x}$ only stretches or
compresses its input, and by how much.

When $\mat{M}$ is symmetric, as is the case for the Laplacian of an
undirected graph, the \emph{spectral theorem} guarantees three useful
facts: all eigenvalues are real; eigenvectors associated with distinct
eigenvalues are orthogonal; and $\mat{M}$ can be written as
$\mat{M} = \mat{Q}\mat{\Lambda}\mat{Q}\transpose$, where $\mat{Q}$ is
orthogonal and $\mat{\Lambda}$ is diagonal.

\paragraph{Worked example.}
Let
\[
\mat{M} = \begin{bmatrix} 3 & 1 \\ 1 & 3 \end{bmatrix}.
\]
Its characteristic polynomial is
$\det(\mat{M} - \lambda \mat{I}) = (3 - \lambda)^2 - 1
= \lambda^2 - 6\lambda + 8 = (\lambda - 2)(\lambda - 4)$, so the
eigenvalues are $\lambda_1 = 2$ and $\lambda_2 = 4$. Solving
$(\mat{M} - 4\mat{I})\vect{v} = \vect{0}$ yields
$\vect{v}_2 = (1,\, 1)\transpose$, and solving
$(\mat{M} - 2\mat{I})\vect{v} = \vect{0}$ yields
$\vect{v}_1 = (1,\, -1)\transpose$. The two eigenvectors are orthogonal,
as the spectral theorem predicts.

Returning to networks, this machinery is exactly what lets us read the
structure of a piping layout from the spectrum of its Laplacian
$\mat{L}$. The smallest eigenvalue is always zero, and the associated
eigenvector $\vect{1}$ says ``every node looks the same from a purely
topological point of view''. The second eigenvalue $\lambda_2$ (the
Fiedler value) measures how tightly the network is connected; its
eigenvector, the \emph{Fiedler vector}, splits the nodes into the two
groups that are hardest to keep together. In a refinery, that split is
a direct hint at where bottlenecks sit and which pipes deserve
redundant backups.

% --------------------------------------------------------------------
\subsection{Cost and flow as a linear model}
\label{sec:cost}

Finally, we put a price on the flow. Associate with each pipe $e$ a
non-negative unit cost $c_e$, collected into a cost vector
$\vect{c} \in \RR^{m}$. The total cost of operating the network at
flow pattern $\vect{f}$ is simply
\begin{equation}
  C(\vect{f}) \;=\; \vect{c}\transpose \vect{f} \;=\;
  \sum_{e \in E} c_e f_e.
  \label{eq:cost}
\end{equation}
Because $C$ is linear in $\vect{f}$, the entire design problem fits in
the familiar template of a \emph{linear program}:
\begin{equation}
  \begin{aligned}
    \text{minimize} \quad & \vect{c}\transpose \vect{f} \\
    \text{subject to} \quad & \mat{B}\vect{f} = \vect{s}, \\
    & \vect{0} \le \vect{f} \le \vect{f}_{\max}.
  \end{aligned}
  \label{eq:lp}
\end{equation}
The equality $\mat{B}\vect{f} = \vect{s}$ is mass conservation from
\eqref{eq:balance}. The inequality $\vect{0} \le \vect{f} \le
\vect{f}_{\max}$ enforces direction and pipe capacity: flow is
non-negative in the chosen orientation, and no pipe exceeds its maximum
throughput.

Formulation \eqref{eq:lp} already suggests the design question we will
answer in the Application section. Given two candidate layouts of the
same refinery (different graph $G$, possibly different capacities
$\vect{f}_{\max}$, and different unit costs $\vect{c}$), the better
layout is the one whose linear program is feasible and whose optimal
value $\vect{c}\transpose \vect{f}^{*}$ is smaller. Eigenvalue analysis
from Section~\ref{sec:eigen} acts as a second, qualitative filter: even
if two layouts produce similar optimal costs, the layout with a larger
Fiedler value is structurally more robust and should be preferred when
failure modes matter.

In Section~3 we apply this framework to a specific refinery piping
network: we construct its adjacency, degree, incidence, and Laplacian
matrices, solve the conservation system by Gaussian elimination,
analyze the Laplacian spectrum, and compare two candidate layouts via
\eqref{eq:lp}.

% =====================================================================
%  3. RESEARCH AND APPLICATION  (target ~3 pages)
% =====================================================================
\section{Research and Application}
\label{sec:application}

% --------------------------------------------------------------------
\subsection{A simplified refinery piping network}
\label{sec:refinery-model}

We model a simplified oil refinery as a directed graph $G = (V, E)$
with six nodes and seven pipes. This is deliberately small enough to
do every calculation by hand and then double-check numerically in
Annex~\ref{annex:python}, while still containing every feature that
makes refinery design interesting: a source, two processing stages, an
internal junction, and two competing output paths.

\paragraph{Nodes.}
\begin{center}
\begin{tabular}{@{}llc@{}}
\toprule
Node & Role & Supply $s_i$ \\
\midrule
$1$ & Crude oil intake (source)   & $+100$ \\
$2$ & Distillation Tower~1 (DT1)  & $0$    \\
$3$ & Distillation Tower~2 (DT2)  & $0$    \\
$4$ & Internal Junction           & $0$    \\
$5$ & Gasoline Outlet (sink)      & $-60$  \\
$6$ & Diesel Outlet (sink)        & $-40$  \\
\bottomrule
\end{tabular}
\end{center}
The external supply vector is therefore
\[
  \vect{s} = (100,\, 0,\, 0,\, 0,\, -60,\, -40)\transpose.
\]
The signs follow Section~\ref{sec:conservation}: the crude intake
produces material, the two outlets consume it, and the internal nodes
are in balance. The total supply is zero ($100 - 60 - 40 = 0$), as it
must be for a closed mass balance.

\paragraph{Pipes.}
Each pipe is a directed edge. The numbering fixes the columns of the
incidence matrix and the entries of the flow vector $\vect{f}$.
\begin{center}
\begin{tabular}{@{}lll@{}}
\toprule
Edge & From $\to$ To & Physical meaning \\
\midrule
$e_1$ & $1 \to 2$ & Crude feed into DT1 \\
$e_2$ & $2 \to 3$ & DT1 heavy fraction into DT2 \\
$e_3$ & $2 \to 4$ & DT1 light fraction into the junction \\
$e_4$ & $3 \to 5$ & DT2 light fraction to gasoline \\
$e_5$ & $3 \to 6$ & DT2 heavy fraction to diesel \\
$e_6$ & $4 \to 5$ & Junction to gasoline \\
$e_7$ & $4 \to 6$ & Junction to diesel \\
\bottomrule
\end{tabular}
\end{center}
% FIGURE: teammate to insert a labelled diagram of this network here.
% Suggested layout:  node 1 on the far left, nodes 2 in the centre-left,
% nodes 3 and 4 stacked in the centre, nodes 5 and 6 stacked on the
% right.  Arrows labelled e1..e7 as per the table above.
We call this the \emph{reference layout} (Layout~A). In
Section~\ref{sec:layout-comparison} we will compare it against a second
layout (Layout~B) in which one pipe has been removed.

% --------------------------------------------------------------------
\subsection{Adjacency, degree, incidence, and Laplacian matrices}
\label{sec:refinery-matrices}

Using the definitions of Section~\ref{sec:graphs}, the four matrices
associated with Layout~A are
\[
\mat{A} =
\begin{bmatrix}
 0 & 1 & 0 & 0 & 0 & 0 \\
 1 & 0 & 1 & 1 & 0 & 0 \\
 0 & 1 & 0 & 0 & 1 & 1 \\
 0 & 1 & 0 & 0 & 1 & 1 \\
 0 & 0 & 1 & 1 & 0 & 0 \\
 0 & 0 & 1 & 1 & 0 & 0
\end{bmatrix},
\qquad
\mat{D} =
\begin{bmatrix}
 1 & 0 & 0 & 0 & 0 & 0 \\
 0 & 3 & 0 & 0 & 0 & 0 \\
 0 & 0 & 3 & 0 & 0 & 0 \\
 0 & 0 & 0 & 3 & 0 & 0 \\
 0 & 0 & 0 & 0 & 2 & 0 \\
 0 & 0 & 0 & 0 & 0 & 2
\end{bmatrix},
\]
\[
\mat{L} = \mat{D} - \mat{A} =
\begin{bmatrix}
 \phantom{-}1 & -1 & \phantom{-}0 & \phantom{-}0 & \phantom{-}0 & \phantom{-}0 \\
 -1 & \phantom{-}3 & -1 & -1 & \phantom{-}0 & \phantom{-}0 \\
 \phantom{-}0 & -1 & \phantom{-}3 & \phantom{-}0 & -1 & -1 \\
 \phantom{-}0 & -1 & \phantom{-}0 & \phantom{-}3 & -1 & -1 \\
 \phantom{-}0 & \phantom{-}0 & -1 & -1 & \phantom{-}2 & \phantom{-}0 \\
 \phantom{-}0 & \phantom{-}0 & -1 & -1 & \phantom{-}0 & \phantom{-}2
\end{bmatrix},
\]
\[
\mat{B} =
\begin{bmatrix}
 \phantom{-}1 & \phantom{-}0 & \phantom{-}0 & \phantom{-}0 & \phantom{-}0 & \phantom{-}0 & \phantom{-}0 \\
 -1 & \phantom{-}1 & \phantom{-}1 & \phantom{-}0 & \phantom{-}0 & \phantom{-}0 & \phantom{-}0 \\
 \phantom{-}0 & -1 & \phantom{-}0 & \phantom{-}1 & \phantom{-}1 & \phantom{-}0 & \phantom{-}0 \\
 \phantom{-}0 & \phantom{-}0 & -1 & \phantom{-}0 & \phantom{-}0 & \phantom{-}1 & \phantom{-}1 \\
 \phantom{-}0 & \phantom{-}0 & \phantom{-}0 & -1 & \phantom{-}0 & -1 & \phantom{-}0 \\
 \phantom{-}0 & \phantom{-}0 & \phantom{-}0 & \phantom{-}0 & -1 & \phantom{-}0 & -1
\end{bmatrix}.
\]
Three quick sanity checks confirm the construction. Every column of
$\mat{B}$ contains exactly one $+1$ (tail) and one $-1$ (head), as it
must for a well-defined directed edge. Every row of $\mat{L}$ sums to
zero. The diagonal of $\mat{D}$ gives the node degrees
$(1, 3, 3, 3, 2, 2)$, which sum to $14 = 2m = 2 \cdot 7$, consistent
with the handshake lemma.

% --------------------------------------------------------------------
\subsection{Conservation equations and Gaussian elimination}
\label{sec:refinery-gauss}

Writing the system $\mat{B}\vect{f} = \vect{s}$ row by row gives the
six node-balance equations
\begin{equation}
\begin{aligned}
  f_1 &= 100,\\
 -f_1 + f_2 + f_3 &= 0,\\
 -f_2 + f_4 + f_5 &= 0,\\
 -f_3 + f_6 + f_7 &= 0,\\
 -f_4 - f_6 &= -60,\\
 -f_5 - f_7 &= -40.
\end{aligned}
\label{eq:refinery-balance}
\end{equation}
There are seven unknowns and six equations, and because the rows of
$\mat{B}$ sum to zero, only five of them are independent. We therefore
expect a two-parameter family of feasible flows.

To find that family explicitly, apply Gaussian elimination to the
augmented matrix $[\mat{B} \mid \vect{s}]$. Working top to bottom with
the pivot in row $k$ used to clear row $k+1$:
\[
\left[\begin{array}{ccccccc|c}
 \phantom{-}1 & \phantom{-}0 & \phantom{-}0 & \phantom{-}0 & \phantom{-}0 & \phantom{-}0 & \phantom{-}0 & \phantom{-}100 \\
 -1 & \phantom{-}1 & \phantom{-}1 & \phantom{-}0 & \phantom{-}0 & \phantom{-}0 & \phantom{-}0 & \phantom{-}0 \\
 \phantom{-}0 & -1 & \phantom{-}0 & \phantom{-}1 & \phantom{-}1 & \phantom{-}0 & \phantom{-}0 & \phantom{-}0 \\
 \phantom{-}0 & \phantom{-}0 & -1 & \phantom{-}0 & \phantom{-}0 & \phantom{-}1 & \phantom{-}1 & \phantom{-}0 \\
 \phantom{-}0 & \phantom{-}0 & \phantom{-}0 & -1 & \phantom{-}0 & -1 & \phantom{-}0 & -60 \\
 \phantom{-}0 & \phantom{-}0 & \phantom{-}0 & \phantom{-}0 & -1 & \phantom{-}0 & -1 & -40
\end{array}\right]
\xrightarrow{\substack{R_2 \gets R_2 + R_1 \\ R_3 \gets R_3 + R_2 \\ R_4 \gets R_4 + R_3 \\ R_5 \gets R_5 + R_4 \\ R_6 \gets R_6 + R_5}}
\]
\[
\left[\begin{array}{ccccccc|c}
 1 & 0 & 0 & 0 & 0 & 0 & 0 & 100 \\
 0 & 1 & 1 & 0 & 0 & 0 & 0 & 100 \\
 0 & 0 & 1 & 1 & 1 & 0 & 0 & 100 \\
 0 & 0 & 0 & 1 & 1 & 1 & 1 & 100 \\
 0 & 0 & 0 & 0 & 1 & 0 & 1 & \phantom{0}40 \\
 0 & 0 & 0 & 0 & 0 & 0 & 0 & \phantom{00}0
\end{array}\right].
\]
The last row becomes $0 = 0$, confirming the rank-$5$ prediction.
Columns $1$--$5$ are pivot columns; columns $6$ and $7$ are free. Back
substitution expresses the \emph{general solution} in terms of the two
free parameters $f_6$ and $f_7$:
\begin{equation}
\begin{aligned}
  f_1 &= 100,\\
  f_2 &= 100 - f_6 - f_7,\\
  f_3 &= f_6 + f_7,\\
  f_4 &= 60 - f_6,\\
  f_5 &= 40 - f_7.
\end{aligned}
\label{eq:general-flow}
\end{equation}
The mass balance by itself therefore tells us that we must push exactly
$100$ units into the intake pipe and that every other flow is
determined once we have chosen how much material to route through the
junction to gasoline ($f_6$) and to diesel ($f_7$). Non-negativity
$\vect{f} \ge \vect{0}$ restricts $(f_6, f_7)$ to the rectangle
$[0, 60] \times [0, 40]$. Choosing the best point in that rectangle is
the optimization problem of Section~\ref{sec:layout-comparison}.

% --------------------------------------------------------------------
\subsection{Laplacian spectrum and network robustness}
\label{sec:refinery-spectrum}

Eigenvalues of the Laplacian $\mat{L}$ tell us whether Layout~A is
merely connected or well-connected. A direct expansion of
$\det(\mat{L} - \lambda \mat{I})$ is painful by hand, but the graph
has two obvious symmetries: swapping nodes $3$ and $4$ leaves it
unchanged (both towers see node $2$ upstream and both outlets
downstream), and swapping nodes $5$ and $6$ does the same. This lets
us split eigenvectors into sign-symmetric and sign-antisymmetric
combinations under each swap, and reduces the $6 \times 6$
eigenproblem into parts we can solve.

\paragraph{Antisymmetric eigenvectors.}
Take $\vect{v} = (0, 0, 0, 0, 1, -1)\transpose$. A direct computation
gives $\mat{L}\vect{v} = (0, 0, 0, 0, 2, -2)\transpose = 2\vect{v}$, so
$\lambda = 2$. Similarly $\vect{v} = (0, 0, 1, -1, 0, 0)\transpose$
gives $\mat{L}\vect{v} = (0, 0, 3, -3, 0, 0)\transpose = 3\vect{v}$, so
$\lambda = 3$.

\paragraph{Symmetric eigenvectors.}
The remaining four eigenvalues come from vectors of the form
$(a, b, c, c, d, d)\transpose$. Substituting into $\mat{L}\vect{v} =
\lambda \vect{v}$ reduces the problem to the $4 \times 4$ non-symmetric
system
\[
\begin{bmatrix}
 1 & -1 & 0 & 0 \\
 -1 & 3 & -2 & 0 \\
 0 & -1 & 3 & -2 \\
 0 & 0 & -2 & 2
\end{bmatrix}
\begin{pmatrix} a \\ b \\ c \\ d \end{pmatrix}
=
\lambda
\begin{pmatrix} a \\ b \\ c \\ d \end{pmatrix}.
\]
Its characteristic polynomial factors as
\[
  \lambda(\lambda - 3)(\lambda^2 - 6\lambda + 4) = 0,
\]
yielding $\lambda \in \{0,\, 3 - \sqrt{5},\, 3,\, 3 + \sqrt{5}\}$.

\paragraph{Full spectrum.}
Combining the two sectors, the eigenvalues of $\mat{L}$ for Layout~A
are
\begin{equation}
  \operatorname{spec}(\mat{L})
  \;=\;
  \bigl\{ 0,\; 3 - \sqrt{5},\; 2,\; 3,\; 3,\; 3 + \sqrt{5} \bigr\}.
  \label{eq:spec-A}
\end{equation}
As a cross-check, the sum of the eigenvalues is $0 + (3-\sqrt{5}) + 2
+ 3 + 3 + (3+\sqrt{5}) = 14$, which agrees with the trace of $\mat{L}$
(sum of node degrees). The Fiedler value is
\[
  \lambda_2
  \;=\; 3 - \sqrt{5}
  \;\approx\; 0.764.
\]
Because $\lambda_2 > 0$, Layout~A is connected: no single pipe failure
can split the refinery into two groups of units that cannot exchange
material, because disconnecting it would require more than one edge
removal. The magnitude $0.764$ will be the benchmark we compare against
when we remove a pipe in Section~\ref{sec:layout-comparison}.

% --------------------------------------------------------------------
\subsection{Cost optimization and layout comparison}
\label{sec:layout-comparison}

Associate the following unit costs with the seven pipes, reflecting
that direct product routes are cheap, routes through the junction are
expensive, and intermediate handling adds a small penalty:
\[
  \vect{c} \;=\; (1,\, 2,\, 1,\, 1,\, 2,\, 3,\, 2)\transpose.
\]
Using the general flow solution \eqref{eq:general-flow}, the total cost
becomes a linear function of the two free variables:
\begin{align*}
  \vect{c}\transpose \vect{f}
  &= 1 \cdot 100
    + 2 \cdot (100 - f_6 - f_7)
    + 1 \cdot (f_6 + f_7)
    + 1 \cdot (60 - f_6) \\
  &\quad {} + 2 \cdot (40 - f_7)
    + 3 \cdot f_6
    + 2 \cdot f_7 \\[2pt]
  &= 440 + f_6 - f_7.
\end{align*}
The linear program from \eqref{eq:lp} therefore reduces to
\[
  \min_{(f_6, f_7) \in [0, 60] \times [0, 40]}
  \bigl( 440 + f_6 - f_7 \bigr),
\]
whose optimum sits at the vertex $(f_6, f_7) = (0, 40)$ with value
\[
  \vect{c}\transpose \vect{f}^{*} \;=\; 400.
\]
Substituting back into \eqref{eq:general-flow} gives the optimal flow
\[
  \vect{f}^{*}
  \;=\;
  (100,\, 60,\, 40,\, 60,\, 0,\, 0,\, 40)\transpose.
\]
Interpretation: send $60$ units through the chain
$1 \to 2 \to 3 \to 5$ to satisfy gasoline demand, and $40$ units
through $1 \to 2 \to 4 \to 6$ to satisfy diesel demand. The two
expensive edges, $e_5$ (DT2$\to$Diesel) and $e_6$
(Junction$\to$Gasoline), carry no flow at optimum; they exist purely as
backup routes.

\paragraph{Layout~B: a pipe is removed.}
Suppose the engineering team considers a cheaper construction in which
pipe $e_4$ (DT2$\to$Gasoline) is not built. Mass balance now forces
$f_6 = 60$ and $f_4$ is eliminated as a variable; one parameter is
left. Minimizing the cost over that one parameter (details in
Annex~\ref{annex:python}) gives
\[
  \vect{c}\transpose \vect{f}^{*}_{\mathrm{B}} \;=\; 460.
\]
That is $60$ units per period more expensive than Layout~A. Using the
same reduction as above (the swap $(1 \leftrightarrow 5)(2
\leftrightarrow 4)(3 \leftrightarrow 6)$ is still an automorphism of
Layout~B), the Laplacian spectrum of Layout~B splits into a symmetric
sector with eigenvalues $\{0,\, 1,\, 3\}$ and an antisymmetric sector
whose eigenvalues are the roots of $\lambda^3 - 8\lambda^2 + 17\lambda
- 8 = 0$, numerically approximately $0.657$, $2.529$, and $4.814$. The
Fiedler value is therefore
\[
  \lambda_2(\mat{L}_{\mathrm{B}}) \;\approx\; 0.657,
\]
a drop of roughly $14\%$ from Layout~A's $0.764$.

\paragraph{Verdict.}
Layout~A is strictly better than Layout~B on both axes.
\begin{center}
\begin{tabular}{@{}lcc@{}}
\toprule
           & Layout~A & Layout~B \\
\midrule
Optimal operating cost $\vect{c}\transpose\vect{f}^{*}$ & $400$   & $460$   \\
Fiedler value $\lambda_2$ (algebraic connectivity)      & $0.764$ & $0.657$ \\
Pipes                                                   & $7$     & $6$     \\
\bottomrule
\end{tabular}
\end{center}
The interesting point is not just that Layout~A wins, but that the two
linear-algebra diagnostics tell consistent stories. The Fiedler value
flags Layout~B as structurally weaker \emph{without any reference to
cost}, and the linear program quantifies the additional cash penalty.
A refinery designer can therefore use the Laplacian spectrum as a
cheap early-stage filter and reserve the full linear-program solve for
the candidates that clear that filter.

% =====================================================================
%  4. CONCLUSION  (target ~0.5 pages)
% =====================================================================
\section{Conclusion}
\label{sec:conclusion}

This report used only the three topics assigned to our group,
\emph{Matrices}, \emph{Gauss Elimination}, and \emph{Eigenvalues}, to
give a complete linear-algebra answer to a genuine industrial design
question: which layout of an oil refinery's piping network is better.
We modelled a simplified refinery as a six-node directed graph,
represented it by its adjacency, degree, incidence, and Laplacian
matrices, solved its mass-balance system by Gaussian elimination to
obtain an explicit two-parameter family of feasible flows, computed
the Laplacian eigenvalues $\{0, 3-\sqrt{5}, 2, 3, 3, 3+\sqrt{5}\}$,
and combined a linear cost model with the balance equations into a
small linear program whose optimum sits at a vertex of the feasible
rectangle.

The three course topics map onto three distinct steps of the
workflow and together cover it end to end. \emph{Matrices} encode the
network: $\mat{A}$ asks ``who is connected to whom'', $\mat{D}$ asks
``how busy is each node'', and $\mat{B}$ translates pipe flows into
node-level balances. \emph{Gauss elimination} solves the resulting
linear system and identifies the free parameters that parametrize the
design space. \emph{Eigenvalues} of the Laplacian provide a single
scalar, the Fiedler value, that measures how robust the layout is to
the loss of a pipe. Comparing Layout~A against Layout~B showed that
these two diagnostics agree in a useful way: the layout with the
larger Fiedler value was also the cheaper one, which means the
eigenvalue analysis can act as an early filter before committing to a
full linear-program solve.

Three directions generalize the study naturally. Realistic refineries
have dozens of nodes rather than six, which makes sparse matrix
techniques and iterative solvers essential but does not change the
underlying linear structure. Pumping and heating costs are nonlinear
in practice, which turns the linear program into a convex one and
replaces Gaussian elimination with Newton-type methods on the KKT
system. Finally, pipe failures are stochastic, and a full treatment
would weight the Fiedler value by an expected-cost model. All three
extensions sit on top of the same matrix--eigenvalue--linear-system
foundation we used here; they are beyond the scope of this course
project but are a direct continuation of it.

% =====================================================================
%  BIBLIOGRAPHY STUB
%  Seeded with the six textbooks the course guidelines recommend.
%  The references teammate should extend this with primary sources
%  (papers, engineering handbooks, refinery design references) as the
%  Application section takes shape.
% =====================================================================
\begin{thebibliography}{9}

\bibitem{anton}
Howard Anton,
\emph{Elementary Linear Algebra (Applications Version)},
Wiley.

\bibitem{lay}
David C. Lay et al.,
\emph{Linear Algebra and Its Applications},
Pearson.

\bibitem{poole}
David Poole,
\emph{Linear Algebra: A Modern Introduction},
Cengage Learning.

\bibitem{strang}
Gilbert Strang,
\emph{Linear Algebra for Everyone},
Wellesley-Cambridge Press.

\bibitem{singh}
Kuldeep Singh,
\emph{Linear Algebra Step by Step},
Oxford University Press.

\bibitem{boyd}
Stephen Boyd and Lieven Vandenberghe,
\emph{Introduction to Applied Linear Algebra: Vectors, Matrices, and
Least Squares},
Cambridge University Press.

% TODO (references teammate): add engineering references on refinery
% piping design, graph theory applied to flow networks, and any primary
% sources used in the Application section.

\end{thebibliography}

% =====================================================================
%  ANNEX A  --  Python verification code
%  Reproduces every numeric claim in Section 3 using numpy and
%  scipy.optimize.linprog.  Requires only numpy and scipy; runs in
%  about half a second on any recent Python.
% =====================================================================
\appendix
\section{Python code for Layouts A and B}
\label{annex:python}

The following Python script verifies the two numerical claims that
drive Section~\ref{sec:application}: the Laplacian spectrum given in
\eqref{eq:spec-A} and the optimal operating costs $\vect{c}\transpose
\vect{f}^{*} = 400$ for Layout~A and $460$ for Layout~B. It also
prints the Fiedler value of each layout and the optimal flow vector
that the linear program returns.

To run the script, save the listing as \texttt{annex\_code.py} and
execute \texttt{pip install numpy scipy} followed by \texttt{python
annex\_code.py}. The output should match the numbers quoted in
Sections~\ref{sec:refinery-spectrum} and~\ref{sec:layout-comparison}.

\begin{verbatim}
import numpy as np
from scipy.optimize import linprog

# --------------------------------------------------------------
# Layout A (reference): 6 nodes, 7 pipes
# Nodes (0-indexed) : 0 = Crude Intake, 1 = DT1, 2 = DT2,
#                     3 = Junction,    4 = Gasoline, 5 = Diesel
# Edges e1..e7      : the directed pipes defined in Section 3.1.
# --------------------------------------------------------------
n = 6
edges_A = [(0, 1),  # e1: Intake -> DT1
           (1, 2),  # e2: DT1 -> DT2
           (1, 3),  # e3: DT1 -> Junction
           (2, 4),  # e4: DT2 -> Gasoline
           (2, 5),  # e5: DT2 -> Diesel
           (3, 4),  # e6: Junction -> Gasoline
           (3, 5)]  # e7: Junction -> Diesel
m_A = len(edges_A)

# Incidence matrix B (6 x 7)
B_A = np.zeros((n, m_A))
for j, (u, v) in enumerate(edges_A):
    B_A[u, j] = +1.0       # tail (flow leaves u)
    B_A[v, j] = -1.0       # head (flow enters v)

# External supply vector s
s = np.array([100, 0, 0, 0, -60, -40], dtype=float)

# Undirected adjacency, degree, and Laplacian
A_A = np.zeros((n, n))
for u, v in edges_A:
    A_A[u, v] = 1.0
    A_A[v, u] = 1.0
D_A = np.diag(A_A.sum(axis=1))
L_A = D_A - A_A

print("Layout A : Laplacian eigenvalues =",
      np.round(np.sort(np.linalg.eigvalsh(L_A)), 4))
# Expected: [0.0, 0.7639, 2.0, 3.0, 3.0, 5.2361]

# Unit cost vector c for Layout A
c_A = np.array([1, 2, 1, 1, 2, 3, 2], dtype=float)

# Linear program:  min c^T f  s.t.  B f = s,  f >= 0
res_A = linprog(c_A, A_eq=B_A, b_eq=s,
                bounds=[(0, None)] * m_A, method='highs')
print("Layout A : optimal flow  =", np.round(res_A.x, 4))
print("Layout A : optimal cost  =", round(res_A.fun, 4))
# Expected flow: [100, 60, 40, 60, 0, 0, 40];  cost: 400.0

# --------------------------------------------------------------
# Layout B : pipe e4 (DT2 -> Gasoline) is removed.
# --------------------------------------------------------------
edges_B = [(0, 1), (1, 2), (1, 3),
           (2, 5), (3, 4), (3, 5)]     # e1, e2, e3, e5, e6, e7
m_B = len(edges_B)

B_B = np.zeros((n, m_B))
for j, (u, v) in enumerate(edges_B):
    B_B[u, j] = +1.0
    B_B[v, j] = -1.0

A_B = np.zeros((n, n))
for u, v in edges_B:
    A_B[u, v] = 1.0
    A_B[v, u] = 1.0
D_B = np.diag(A_B.sum(axis=1))
L_B = D_B - A_B

print("Layout B : Laplacian eigenvalues =",
      np.round(np.sort(np.linalg.eigvalsh(L_B)), 4))
# Expected: [0.0, 0.6571, 1.0, 2.5293, 3.0, 4.8136]

# Cost vector without c4 : (c1, c2, c3, c5, c6, c7)
c_B = np.array([1, 2, 1, 2, 3, 2], dtype=float)

res_B = linprog(c_B, A_eq=B_B, b_eq=s,
                bounds=[(0, None)] * m_B, method='highs')
print("Layout B : optimal flow  =", np.round(res_B.x, 4))
print("Layout B : optimal cost  =", round(res_B.fun, 4))
# Expected flow: [100, 0, 100, 0, 60, 40];  cost: 460.0

# --------------------------------------------------------------
# Fiedler values (algebraic connectivity) for the two layouts
# --------------------------------------------------------------
fiedler_A = np.sort(np.linalg.eigvalsh(L_A))[1]
fiedler_B = np.sort(np.linalg.eigvalsh(L_B))[1]
print(f"Fiedler value Layout A = {fiedler_A:.4f}")
print(f"Fiedler value Layout B = {fiedler_B:.4f}")
# Expected: Layout A ~ 0.7639,  Layout B ~ 0.6571
\end{verbatim}

Running the script reproduces every number quoted in
Section~\ref{sec:application}: the Laplacian spectrum
$\{0,\, 3-\sqrt{5},\, 2,\, 3,\, 3,\, 3+\sqrt{5}\}$ for Layout~A and the
numerically-approximated spectrum $\{0,\, 0.657,\, 1,\, 2.529,\, 3,\,
4.814\}$ for Layout~B, the optimal flows
$\vect{f}^{*}_{\mathrm{A}} = (100, 60, 40, 60, 0, 0, 40)\transpose$ and
$\vect{f}^{*}_{\mathrm{B}} = (100, 0, 100, 0, 60, 40)\transpose$, and
the optimal operating costs $400$ and $460$. The script is short
enough to be pasted into a Jupyter notebook, and it is easy to extend:
adding a new pipe is one more entry in \texttt{edges\_A}, and a new
cost experiment is one more run of \texttt{linprog}.

\end{document}
